\clearpage\begingroup\SourceLeaf{373b}\ArithmeticHead{42}{L}{STÄCKEL, GAUSS ALS GEOMETER.}
\renewcommand{\thefootnote}{\arabic{footnote})}\setcounter{footnote}{0}
\begin{center}
\textls[100]{C. Sonstige Beiträge zur Axiomatik.}\par\medskip
\textbf{10.}\par\smallskip
\textls[100]{Weitere Untersuchungen über die Grundlagen der Geometrie.}
\end{center}
\noindent\hspace*{\parindent}Wenn von den Untersuchungen die Rede ist, die \textsc{Gauss} über die Grundlagen
der Geometrie angestellt hat, so denkt man dabei vor allem an die
Entdeckung der nichteuklidischen Geometrie. \textsc{Gauss} hat sich jedoch keineswegs
auf das Parallelenaxiom beschränkt, er hat sich vielmehr noch mit einer
Reihe anderer Fragen beschäftigt, die man heute ebenfalls der Axiomatik zuweisen
würde. Hierüber soll zum Schluss dieses Abschnittes berichtet werden.
\par
Es kann nicht Wunder nehmen, dass die üblichen Darstellungen der
euklidischen Geometrie einen Mann, der an die Schärfe der Begriffsbestimmungen
und die Strenge der Ableitungen hohe Forderungen stellte, in mehr als einem
Punkte nicht befriedigten. Sein tiefdringender Blick erkannte hier Lücken,
die zum Teil erst nach Jahrzehnten von anderen Geometern aufgedeckt worden
sind. Zum Beispiel spricht \textsc{Gauss} in dem Brief an \textsc{Bolyai} vom 6. März 1832
von dem »Teil des Planums, der zwischen drei Geraden liegt« und macht dazu
die Anmerkung: »Bei einer vollständigen Durchführung müssen solche Worte
\clearpage\SourceLeaf{374}\ArithmeticHead{43}{R}{DIE GRUNDLAGEN DER GEOMETRIE.}\setcounter{footnote}{0}
\noindent wie \textls[100]{zwischen} auch erst auf klare Begriffe gebracht werden, was sehr gut
angeht, was ich aber nirgends geleistet finde« (W. VIII, S. 222).
\par
Die Erklärung der geraden Linie war Gegenstand des Gespräches gewesen,
das \textsc{Gauss} und \textsc{Wolfgang Bolyai} bei ihrem ersten gemeinsamen
Spaziergange im Herbst 1796 geführt hatten. Wie \textsc{Johann Bolyai} erzählt,
erwiderte \textsc{Gauss} auf die Äusserungen \textsc{Wolfgangs}: »Ja wahrlich, die Gerade
wird schändlich behandelt; sie ist in der Tat die Linie, welche sich in sich
selbst dreht« (\textsc{Bol.} S. 197). Dieselbe Erklärung hat \textsc{Gauss} in einer Vorlesung
über praktische Astronomie gegeben, die \textsc{Lübsen} im Jahre 1830 bei ihm gehört
hat (W. VIII, S. 196); auch die weitere Bemerkung bei \textsc{Lübsen}, das
angegebene Merkmal sei praktisch wichtig, z. B. bei der Justierung eines
Fernrohres, bei der richtigen Bohrung eines Zylinders usw., ist wohl \textsc{Gauss}schen
Ursprungs.
\par
Im \textit{Tagebuch} steht unter dem 28. Juli 1798 die Eintragung: »Plani possibilitatem
demonstravi« (T. Nr. 72). Was \textsc{Gauss} hiermit meinte, zeigt eine
Stelle in dem Briefe an \textsc{Bessel} vom 27. Januar 1829, die Erklärung der Ebene
als einer Fläche, in der die irgend zwei Punkte verbindende gerade Linie
ganz liegt, enthalte mehr, als zur Bestimmung der Fläche nötig ist, und involviere
tacite ein Theorem, das erst bewiesen werden müsse (W. VIII, S. 200);
in ähnlicher Weise äussert sich \textsc{Gauss} auch in einer wohl aus der gleichen
Zeit stammenden Aufzeichnung (W. VIII, S. 194). Auch in dem Brief an
W. \textsc{Bolyai} vom 6. März 1832 erklärt es \textsc{Gauss} für unerlässlich, »die Möglichkeit
eines Planums zu erweisen« (W. VIII, S. 224). Ein solcher Beweis steht
im Handbuch 19 Be, S. 153 (W. VIII, S. 194); durch die unmittelbar vorhergehenden
Notizen, die den Dreiecks-Inhalt und das Tetraeder-Volumen in der
nichteuklidischen Geometrie betreffen (W. VIII, S. 226—228), ist als Zeit der
Niederschrift der März 1832 gesichert. Die Ebene denkt sich \textsc{Gauss} erzeugt
durch die Drehung des einen Schenkels eines rechten Winkels um den anderen,
festgehaltenen Schenkel. Auf Anregungen von \textsc{Gauss} gehen wohl auch die
Abhandlungen von \textsc{Deahna} (1837) und \textsc{Gerling} (1840) über die Erklärung
der Ebene zurück\footnote{\textsc{Deahna}, \textit{Demonstratio theorematis esse superficiem planam}, Marburg 1837; \textsc{Chr. L. Gerling}, \textit{Fragment über die Begründung des Begriffs der Ebene}, \textsc{Crelles} Journal, Bd. 20, 1840, S. 332. \textsc{Baltzer} bemerkt in der zweiten Auflage seiner \textit{Elemente}, Bd. II, 1867, § 4, \textsc{Gauss} sei der Meinung gewesen, \textsc{Deahnas} Darstellung lasse sich von einigen Mängeln, die in ihr anzutreffen seien, befreien; vgl. auch W. \textsc{Killing}, \textit{Einführung in die Grundlagen der Geometrie}, Bd. II, Paderborn 1898, S. 183.}.
\par\vfill\hfill{\footnotesize 6*}\hspace*{2em}
\clearpage\SourceLeaf{375}\ArithmeticHead{44}{L}{STÄCKEL, GAUSS ALS GEOMETER.}\setcounter{footnote}{0}
\noindent\hspace*{\parindent}Bei der Lehre von den Vielecken pflegt man stillschweigend oder ausdrücklich
voraussetzen, dass der Umfang sich selbst nicht schneidet. \textsc{Gauss}
hat schon früh die Frage ins Auge gefasst, was man unter dem Inhalt eines
beliebigen Vielecks zu verstehen habe; in Nr. 24 dieses Aufsatzes wird man
hierüber Genaueres finden. Bei seinen Untersuchungen über die allgemeine
Lehre von den krummen Flächen ist er auf den Gegenstand zurückgekommen
und hat beliebige Figuren betrachtet, deren Umfang sich selbst schneidet
(Brief an \textsc{Olbers} vom 20. Oktober 1825, Br. G.-O. 2, S. 431, W. VIII, S. 399;
vgl. auch W. IV, S. 227). Auch die Zerlegung der Vielecke in Dreiecke hat
er untersucht (W. VIII, S. 280); sein Verfahren führt zu einer Herleitung der
Winkelsumme des \(n\)-Ecks, die dem üblichen, unzulänglichen Induktionsbeweise
vorzuziehen ist.
\par
Eine Anfrage \textsc{Gerlings} vom 20. Juni 1846 über die Unterscheidung rechts-
und linksgewundener Schrauben (W. VIII, S. 247) veranlasste \textsc{Gauss} zu Ausführungen
über die Begriffe rechts und links, die er »ein Kernstück eines viel
ausgedehntern Systems« nennt (Brief vom 23. Juni 1846, W. VIII, S 249).
Er war bereits in der Selbstanzeige der zweiten Abhandlung über die biquadratischen
Reste vom 15. April 1831 (W. II, S. 177), in der er seine geometrische
Versinnlichung der komplexen Grössen darlegt, auf den Unterschied
von rechts und links eingegangen und hatte bemerkt, dieser Unterschied sei
»sobald man vorwärts und rückwärts in der Ebene und oben und unten in
Beziehung auf die beiden Seiten der Ebene einmal (nach Gefallen) festgesetzt
hat, \textls[100]{in sich} völlig bestimmt, wenn wir gleich unsere Anschauung dieses
Unterschiedes andern \textls[100]{nur} durch Nachweisung an wirklich vorhandenen materiellen
Dingen mitteilen können«. In einer Fussnote hatte er hinzugefügt:
»Beide Bemerkungen hat schon \textsc{Kant} gemacht, aber man begreift nicht, wie
dieser scharfsinnige Philosoph in der ersten einen Beweis für seine Meinung,
dass der Raum \textls[100]{nur} Form unserer äusseren Anschauung sei, zu finden glauben
konnte, da die zweite so klar das Gegenteil und dass der Raum unabhängig
von unserer Anschauungsart eine reelle Bedeutung haben muss, beweiset«
(vgl. auch W. X\textsubscript{1}, S. 409); eine ähnliche Bemerkung enthält der Brief an
\clearpage\SourceLeaf{376}\ArithmeticHead{45}{R}{DIE GRUNDLAGEN DER GEOMETRIE.}\setcounter{footnote}{0}
\noindent \textsc{Schumacher} vom 8. Februar 1846 (W. VIII, S. 247)\footnote{Vgl. noch E. \textsc{Study}, \textit{Die Begriffe Links, Rechts, Windungssinn und Drehungssinn}, Archiv der Mathematik und Physik, 3. Reihe, Bd. 21, 1913, S. 193; hier wird auf den Briefwechsel zwischen \textsc{Gauss} und \textsc{Gerling} ausführlich Bezug genommen.}. Einen zweiten Grund
gegen \textsc{Kants} Meinung hat \textsc{Gauss} in dem Brief an W. \textsc{Bolyai} vom 6. März
1832 vorgebracht. »Gerade in der Unmöglichkeit zwischen \(\Sigma\) [Euklidischer
Geometrie] und \(S\) [nichteuklidischer Geometrie] a priori zu entscheiden, liegt
der klarste Beweis, dass \textsc{Kant} Unrecht hatte zu behaupten, der Raum sei nur
\textls[100]{Form} unserer Anschauung« (W. VIII, S. 224).
\par
In dasselbe Kapitel wie die Erörterungen über die Begriffe von Rechts
und Links gehören die Einführung der gerichteten geraden Linien (W. VIII,
S. 408), die Unterscheidung zwischen den beiden zu einem grössten Kreise
der Kugel gehörenden Polen (W. VII, S. 177, IV, S. 221) und die Sätze, dass
symmetrische sphärische Dreiecke flächengleich, symmetrische Raumstücke
volumengleich sind. \textsc{Gerling} hatte den ersten Satz durch Zerlegung in Teil-
Dreiecke bewiesen, die paarweise kongruent sind (Brief an \textsc{Gauss} vom 25. März
1813, W. VIII, S. 240). Als er am 26. Februar 1844 darauf zurückkam,
forderte ihn \textsc{Gauss} auf, den zweiten zu beweisen (Brief vom 8. April 1844,
W. VIII, S. 241). \textsc{Gerling} konnte auch hier zeigen, dass die Gebilde sich
in Pyramiden zerlegen lassen, die paarweise kongruent sind (Brief vom 15. April
1844, W. VIII, S. 242)\footnote{\textsc{Gerlings} Beweis ist von \textsc{Hessel} vereinfacht worden: \textit{Einige neue Beweise von Lehrsätzen aus der Elementar-Stereometrie}, Archiv der Mathematik und Physik, 1. Reihe, Bd. 7, 1846, S. 284; \textsc{Hessel} bemerkt, dass \textsc{Gerling} durch \textsc{Gauss} zu seinen Untersuchungen veranlasst worden sei.}. Nunmehr warf \textsc{Gauss} die Frage auf, ob man in
ähnlicher Weise, unabhängig von der Exhaustionsmethode, zeigen könne, dass
Pyramiden von gleicher Grundfläche und gleicher Höhe gleichen Rauminhalt
haben (Brief vom 17. April 1844, W. VIII, S. 244)\footnote{Es ist nicht ausgeschlossen, dass \textsc{Gauss} diese Fragestellung dem \textit{Tentamen} \textsc{Wolfgang Bolyais} verdankte; dieser hatte die Frage von der »endlichen Gleichheit« bei Flächenstücken ausführlich untersucht und dazu bemerkt, ob eine beliebige dreiseitige Pyramide durch endliche Gleichheit auf ein Prisma zurückgeführt werden könne oder nicht, sei noch nicht klargestellt (\textit{Tentamen}, t. II, S. 175, ed. secunda, Budapest 1904, S. 241; vgl. \textsc{Bol.} S. 40 und 188).}, aber hier gelangte \textsc{Gerling}
nicht zum Ziele (Brief vom 7. Juli 1844, W. VIII, S. 245). Durch die
Herausgabe der bis dahin unbekannten Briefe von \textsc{Gauss} und \textsc{Gerling} im
achten Bande der Werke (1900) wurde die Aufmerksamkeit auf die Frage
der Volumengleichheit der Polyeder gelenkt, und so könnte man letzten Endes
\clearpage\SourceLeaf{377}\ArithmeticHead{46}{L}{STÄCKEL, GAUSS ALS GEOMETER.}\setcounter{footnote}{0}
\noindent den Beweis \textsc{Dehns}, dass die Exhaustionsmethode bei der Volumenbestimmung
unentbehrlich ist\footnote{M. \textsc{Dehn}, \textit{Über raumgleiche Polyeder}, Göttinger Nachrichten 1900, S. 345; \textit{Über den Rauminhalt}, Math. Annalen, Bd. 55, 1901, S. 465; vgl. jedoch schon R. \textsc{Bricard}, \textit{Sur une question de géométrie rélative aux polyèdres}, Nouv. ann. de math., série 3, t. 15, 1896, S. 331 und G. \textsc{Sforza}, \textit{Un’ osservazione sull’ equivalenza dei poliedri per congruenza delle parti}, Periodico di mat., t. 12, 1897, S. 105.}, auf eine Anregung von \textsc{Gauss} zurückführen.
\par\medskip\begin{center}\rule{0.18\textwidth}{0.3pt}\end{center}
\clearpage\endgroup
